\documentclass[10pt,a4paper]{amsart}
\usepackage[margin=19mm]{geometry}
\usepackage{amsmath,amssymb,mathtools}
\usepackage[scr=rsfs,cal=euler]{mathalfa}
\usepackage{xfrac}
\usepackage[unicode,hidelinks,bookmarks=false]{hyperref}
\newcommand{\ie}{i.e.\ }
\newcommand{\cf}{cf.\ }
\newcommand{\resp}{respectively\ }
\newcommand{\Subsection}{\subsection}
\providecommand{\nobreakdash}{\nobreak\hbox{-}}
\setlength{\emergencystretch}{2em}
\multlinegap0pt
\usepackage{tikz}
\usepackage{amssymb}
\usepackage{bm}
\newtheorem{lem}{Lemma}
\newcommand{\compi}{\textnormal{i}}
\newcommand{\ind}{\mathbf{1}}
\title{The Density Formula Approach for Non-reversible Isomorphism Theorems, with Applications}
\author{Qinghua (Devon) Ding, Venkat Anantharam}
\date{Source: arXiv:2503.12859v2 (CC BY 4.0)}
\begin{document}
\maketitle

\noindent\emph{Source and licence.} The Density Formula Approach for Non-reversible Isomorphism Theorems, with Applications. Version arXiv:2503.12859v2 (CC BY 4.0), distributed by arXiv under the Creative Commons Attribution 4.0 International licence, http://creativecommons.org/licenses/by/4.0/, as recorded in the arXiv licence metadata for this exact version and shown on the abstract page https://arxiv.org/abs/2503.12859v2. Full licence notice: \texttt{licenses/arxiv-2503.12859-CC-BY-4.0.txt}.\par\smallskip

\noindent\textit{Editorial scope.} Counted unit: the article's lem lem:nonneg (source0.tex lines 1725--1728). The article's complete original proof of it is delivered with it (source0.tex lines 1730--1779, one \texttt{proof} environment of 50 lines). No mathematics is written, repaired or re-derived: the statement and the proof are the article's own lines, in the article's own notation, and no retained line is substituted or re-flowed.  Retained uncounted as the source-local context the proof needs: lines 116--118.  Nothing was omitted from any retained span, so the package carries no omission record.  No retained line contains a citation, so the wrapper carries no bibliography and no bibliography entry.  No retained line contains a figure environment, an image-inclusion command or drawing code, so no image is delivered and the package declares no asset dependency.  Editorial formatting only, in addition: an AMS \texttt{amsart} preamble that restates the article's own declarations and macros used by the retained lines, namely the environments \texttt{lem} on separate counters, together with the standard packages those lines need.  The retained environments print in this document as Lemma 1 in order of appearance, whereas the article prints the same items with the numbering of its own full document.  The complete native source of the article as distributed in the arXiv e-print for this exact version is bundled unchanged as \texttt{sources/174885/source0.tex}.
\begin{equation}        \label{eq:diffform}
d \bar{\phi} \wedge d \phi = (du - \compi dv) \wedge (du + \compi dv) = 2 \compi du \wedge dv = \compi dl \wedge d\theta,
\end{equation}

\vspace{.3em} \begin{lem}
\label{lem:nonneg}
Let $G$ be an inverse M-matrix with a positive definite symmetric part. For any nonnegative function $f:\mathbb{R}_+^n\rightarrow \mathbb{R}_+$ with subexponential growth, the integral $\int_{\mathbb{C}^n}\, \phi_a\bar{\phi}_b f(\phi\bar{\phi}) d\mu(\phi)$ is finite and nonnegative for any $a,b\in [n]$.
\end{lem}

\begin{proof}
To see that
    %it
the integral is nonnegative, note that
    \begin{align}   \label{eq:appceqn}
    \int_{\mathbb{C}^n}\, \phi_a\bar{\phi}_b f(\phi\bar{\phi}) d\mu(\phi) & = \frac{1}{(2\pi\compi)^n}\int_{\mathbb{C}^n}\, \phi_a\bar{\phi}_b f(\phi\bar{\phi}) e^{-\sum_{ij} A_{ij}\phi_i\bar{\phi}_j} d\bar{\phi}d\phi \nonumber\\
    & = \frac{1}{(2\pi)^n}\int_{\mathbb{R}_+^n}\int_{[0,2\pi)^n}\, \sqrt{l_a l_b}e^{\compi(\theta_a-\theta_b)} f(l) e^{-\sum_{ij} A_{ij}\sqrt{l_i l_j}e^{\compi(\theta_i-\theta_j)}} d\theta dl \nonumber\\
    & = \int_{\mathbb{R}_+^n}\sqrt{l_a l_b} f(l) e^{-\sum_{i} A_{ii}l_i}\bigg(\frac{1}{(2\pi)^n}\int_{[0,2\pi)^n}\, e^{\compi(\theta_a-\theta_b)}   e^{-\sum_{i\neq j} A_{ij}\sqrt{l_i l_j}e^{\compi(\theta_i-\theta_j)}} d\theta\bigg) dl,
    \end{align}
where we have used equation \eqref{eq:diffform} in writing the second equality.

Moreover, we have
    \begin{align*}
    e^{-\sum_{i\neq j} A_{ij}\sqrt{l_i l_j}e^{\compi(\theta_i-\theta_j)}} & = \prod_{i\neq j}e^{- A_{ij}\sqrt{l_i l_j}e^{\compi(\theta_i-\theta_j)}}\\
    & = \prod_{i\neq j}\sum_{k_{ij}\in\mathbb{N}} \frac{(- A_{ij}\sqrt{l_i l_j}e^{\compi(\theta_i-\theta_j)})^{k_{ij}}}{k_{ij}!}\\
    & = \sum_{\{k_{ij}\in\mathbb{N}, \,i\neq j\}}\prod_{i\neq j} \frac{(- A_{ij})^{k_{ij}}(l_i l_j)^{\frac{1}{2}k_{ij}}}{k_{ij}!}e^{\compi k_{ij}(\theta_i-\theta_j)}.
    \end{align*}

Fix some $l\in\mathbb{R}_+^n$. The inner integral in equation \eqref{eq:appceqn} becomes
    \begin{align*}
    \frac{1}{(2\pi)^n}\int_{[0,2\pi)^n}\, & e^{\compi(\theta_a-\theta_b)}   e^{-\sum_{i\neq j} A_{ij}\sqrt{l_i l_j}e^{\compi(\theta_i-\theta_j)}} d\theta \\
    &= \frac{1}{(2\pi)^n}\int_{[0,2\pi)^n}\, e^{\compi(\theta_a-\theta_b)}\bigg(\sum_{\{k_{ij}\in\mathbb{N}, \,i\neq j\}}\prod_{i\neq j} \frac{(- A_{ij})^{k_{ij}}(l_i l_j)^{\frac{1}{2}k_{ij}}}{k_{ij}!}e^{\compi k_{ij}(\theta_i-\theta_j)}\bigg) d\theta  \\
    &= \frac{1}{(2\pi)^n}\int_{[0,2\pi)^n}\, e^{\compi(\theta_a-\theta_b)}\bigg(\sum_{\{k_{ij}\in\mathbb{N}, \,i\neq j\}}\bigg(\prod_{i\neq j} \frac{(- A_{ij})^{k_{ij}}(l_i l_j)^{\frac{1}{2}k_{ij}}}{k_{ij}!}\bigg)e^{\sum_{i\neq j}\compi k_{ij}(\theta_i-\theta_j)}\bigg) d\theta  \\
    &= \sum_{\{k_{ij}\in\mathbb{N}, \,i\neq j\}}\bigg(\prod_{i\neq j} \frac{(- A_{ij})^{k_{ij}}(l_i l_j)^{\frac{1}{2}k_{ij}}}{k_{ij}!}\bigg)\cdot\bigg(\frac{1}{(2\pi)^n}\int_{[0,2\pi)^n}\, e^{\compi(\theta_a-\theta_b)}e^{\sum_{i\neq j}\compi k_{ij}(\theta_i-\theta_j)} d\theta\bigg).
    \end{align*}

Notice that $\frac{1}{(2\pi)^n}\int_{[0,2\pi)^n}\, e^{\compi(\theta_a-\theta_b)}e^{\sum_{i\neq j}\compi k_{ij}(\theta_i-\theta_j)} d\theta=\frac{1}{(2\pi)^n}\int_{[0,2\pi)^n}\, e^{\sum_{i\neq j}\compi r_i\theta_i} d\theta$ where $r_i:=\sum_{j\neq i} (k_{ij}- k_{ji})+\ind(i=a)-\ind(i=b)$. Since $\frac{1}{2\pi}\int_{[0,2\pi)}e^{\compi r_i\theta_i}d\theta_i=\ind(r_i=0)$, we know that $\frac{1}{(2\pi)^n}\int_{[0,2\pi)^n}\, e^{\compi(\theta_a-\theta_b)}e^{\sum_{i\neq j}\compi k_{ij}(\theta_i-\theta_j)} d\theta=\ind(r_i=0,\forall i\in[n])$. Therefore, we conclude that
    \begin{align*}
    \frac{1}{(2\pi)^n}\int_{[0,2\pi)^n}\, & e^{\compi(\theta_a-\theta_b)}   e^{-\sum_{i\neq j} A_{ij}\sqrt{l_i l_j}e^{\compi(\theta_i-\theta_j)}} d\theta \\
    &= \sum_{\{k_{ij}\in\mathbb{N}, \,i\neq j\}}\bigg(\prod_{i\neq j} \frac{(- A_{ij})^{k_{ij}}(l_i l_j)^{\frac{1}{2}k_{ij}}}{k_{ij}!}\bigg)\cdot \ind(r_i=0,\forall i\in[n])\\
    &= \sum_{\{k_{ij}\in\mathbb{N},\,i\neq j:\, r_i(k)=0,\forall i\in[n]\}}\bigg(\prod_{i\neq j} \frac{(- A_{ij})^{k_{ij}}(l_i l_j)^{\frac{1}{2}k_{ij}}}{k_{ij}!}\bigg).
    \end{align*}

The original integral should be
    \begin{align*}
    \int_{\mathbb{C}^n}\, & \phi_a\bar{\phi}_b f(\phi\bar{\phi}) d\mu(\phi) \\
    & = \int_{\mathbb{R}_+^n}\sqrt{l_a l_b} f(l) e^{-\sum_{i} A_{ii}l_i} \cdot \sum_{\{k_{ij}\in\mathbb{N},\,i\neq j:\, r_i(k)=0,\forall i\in[n]\}}\bigg(\prod_{i\neq j} \frac{(- A_{ij})^{k_{ij}}(l_i l_j)^{\frac{1}{2}k_{ij}}}{k_{ij}!}\bigg) dl\geq 0.
    \end{align*}

Here, we used the sign property of M-matrices, which guarantees $-A_{ij}\geq 0, \forall i\neq j$. We now justify the series expansions and the integration-summation exchanges, critically using positive definiteness. Note that we have
    \begin{align*}
    \int_{\mathbb{R}_+^n}\sqrt{l_a l_b} & f(l) e^{-\sum_{i} A_{ii}l_i}\bigg(\frac{1}{(2\pi)^n}\int_{[0,2\pi)^n}\, \prod_{i\neq j}\sum_{k_{ij}\in\mathbb{N}} \bigg|e^{\compi(\theta_a-\theta_b)}\frac{(- A_{ij}\sqrt{l_i l_j}e^{\compi(\theta_i-\theta_j)})^{k_{ij}}}{k_{ij}!}\bigg| d\theta\bigg) dl\\
    &= \int_{\mathbb{R}_+^n}\sqrt{l_a l_b} f(l) e^{-\sum_{i} A_{ii}l_i}\bigg(\frac{1}{(2\pi)^n}\int_{[0,2\pi)^n}\, \prod_{i\neq j}\sum_{k_{ij}\in\mathbb{N}} \frac{(-A_{ij})^{k_{ij}}(l_i l_j)^{\frac{1}{2}k_{ij}}}{k_{ij}!} d\theta\bigg) dl\\
    &= \int_{\mathbb{R}_+^n}\sqrt{l_a l_b} f(l) e^{-\sum_{i} A_{ii}l_i}\bigg(\frac{1}{(2\pi)^n}\int_{[0,2\pi)^n}\, \prod_{i\neq j}e^{-A_{ij}\sqrt{l_i l_j}} d\theta\bigg) dl\\
    &= \int_{\mathbb{R}_+^n}\sqrt{l_a l_b} f(l) e^{-\sum_{i} A_{ii}l_i} \prod_{i\neq j}e^{-A_{ij}\sqrt{l_i l_j}} dl\\
    &= \int_{\mathbb{R}_+^n}\sqrt{l_a l_b} f(l) e^{-\sum_{i} A_{ii}l_i} e^{-\sum_{i\neq j}A_{ij}\sqrt{l_i l_j}} dl\\
    &= \int_{\mathbb{R}_+^n}\sqrt{l_a l_b} f(l) e^{-\langle \sqrt{l}, A\sqrt{l}\rangle} dl.
    \end{align*}
Since $A$ has positive definite symmetric part, we know there exists $c>0$, such that $\langle \sqrt{l}, A\sqrt{l}\rangle \geq c\|\sqrt{l}\|^2=c\sum_il_i$. Therefore, the above integral should be less than $ \int_{\mathbb{R}_+^n}\sqrt{l_a l_b} f(l) e^{-c\sum_i l_i} dl.$ Now $f(l)$ having subexponential growth guarantees this integral to be finite. Then, the dominated convergence theorem guarantees the expansion and integration-summation exchanges.
\end{proof}

\end{document}
